% ch8_a 宇宙学 (书页 323-335 / p338-p350)
% 第8章 Cosmology
\chapter{宇宙学}

\section{最大对称宇宙}

% ---- 书页 323 / p338 ----
当代的宇宙学模型基于这样一种观念：宇宙在各处都大致相同——这一立场有时被称为\textbf{哥白尼原理}（Copernican principle）。乍一看，这种主张似乎很疯狂；比方说，太阳的中心与星际空间那荒凉的寒冷几乎毫无相似之处。但我们把哥白尼原理理解为仅仅适用在最大的尺度上，在那种尺度上密度的局域变化已被平均掉。它在这种尺度上的有效性体现在许多不同的观测之中，例如对星系的计数以及对弥漫 X 射线和 $\gamma$ 射线背景的观测；不过它最清晰的表现是 3K 的宇宙微波背景（cosmic microwave background，CMB）。虽然我们现在知道，微波背景辐射并非完美光滑（其实从来也没有人指望它完美光滑），但对规则性的偏离只有 $10^{-5}$ 的量级或更小，对于大尺度上时空的近似描述来说，这当然是足够好的基础。

哥白尼原理与流形可能具有的两个在数学上更精确的性质相关：各向同性（isotropy）与均匀性（homogeneity）。\textbf{各向同性}作用于流形中某个特定的点，它断言：无论你朝哪个方向看，空间看起来都一样。更正式地说，如果对于 $T_pM$ 中的任意两个矢量 $V$ 和 $W$，都存在 $M$ 的一个等距变换（isometry），使得 $W$ 在此等距变换下的推前与 $V$ 平行（$V$ 本身不被推前），那么流形 $M$ 就称为绕点 $p$ 各向同性。微波背景的观测所指示的，正是空间的各向同性。

\textbf{均匀性}则是度规在整个流形上都相同的陈述。换句话说，给定 $M$ 中任意两点 $p$ 和 $q$，都存在一个把 $p$ 变到 $q$ 的等距变换。注意，均匀性与各向同性之间没有必然的联系；一个流形可以均匀却无一处各向同性（比如通常度规下的 $\mathbf{R}\times S^2$），也可以绕某一点各向同性却不均匀（比如一个锥面，它绕自己的顶点各向同性，但显然不均匀）。另一方面，如果一个空间\emph{处处}各向同性，那么它就是均匀的。同样，如果它绕某一点各向同性并且还均匀，那么它将绕每一点都各向同性。既然各向同性有着充分的观测证据，而哥白尼原理又要我们相信自己并非宇宙的中心、因而别处的观者也应观测到各向同性，那么从今往后我们就同时假定均匀性和各向同性。

% ---- 书页 324 / p339 ----
均匀性和各向同性的用处在于，它们蕴含空间是最大对称的（maximally symmetric）。把各向同性想成转动下的不变性，把均匀性想成经过适当推广的平移下的不变性。于是均匀性和各向同性合在一起就蕴含：一个空间具有其可能的最大数目的 Killing 矢量。哥白尼原理的一个极端应用，是坚持认为时空本身就是最大对称的。事实上结果将表明并非如此；从观测上我们知道，宇宙在\emph{空间}上是均匀且各向同性的，但在全部\emph{时空}上却不是。不过，从考察最大对称的时空出发是有趣的（它们毕竟只是“仅空间最大对称”这一更普遍情形的特例）。我们将会看到，在某种意义上，这样的宇宙是广义相对论的“基态”。与本章后面的部分相比，这一讨论与观测到的宇宙关系较远；注重经验的读者不妨直接跳到下一节去。

我们在第 3 章中曾提到，带度规 $g_{\mu\nu}$ 的最大对称 $n$ 维流形，其 Riemann 张量可以写成
\begin{equation*}
R_{\rho\sigma\mu\nu} = \kappa\left(g_{\rho\mu}g_{\sigma\nu} - g_{\rho\nu}g_{\sigma\mu}\right),
\tag{8.1}
\end{equation*}
其中 $\kappa$ 是 Ricci 曲率的归一化量度，
\begin{equation*}
\kappa = \frac{R}{n(n-1)},
\tag{8.2}
\end{equation*}
而 Ricci 标量 $R$ 在流形上将是一个常数。由于在任何一个单点上我们总可以把度规化为规范形式（$g_{\mu\nu}=\eta_{\mu\nu}$），最大对称流形的种类在局域上就由度规的号差以及常数 $\kappa$ 的符号来刻画。需要“局域”这个限定词，是为了顾及可能的全局差异，比如平面与环面之间的差异。我们感兴趣的是号差为 $(-+++)$ 的度规。对于曲率为零（$\kappa=0$）的情形，最大对称时空是众所周知的；它恰恰就是 Minkowski 空间，其度规为
\begin{equation*}
\mathrm{d}s^2 = -\mathrm{d}t^2 + \mathrm{d}x^2 + \mathrm{d}y^2 + \mathrm{d}z^2.
\tag{8.3}
\end{equation*}
Minkowski 空间的共形图在附录 H 中导出。

曲率为正（$\kappa>0$）的最大对称时空称为 \textbf{de Sitter 空间}（de Sitter space）。考虑一个五维 Minkowski 空间，度规为 $\mathrm{d}s_5^2 = -\mathrm{d}u^2 + \mathrm{d}x^2 + \mathrm{d}y^2 + \mathrm{d}z^2 + \mathrm{d}w^2$，并在其中嵌入一个由下式给出的双曲面（hyperboloid）
\begin{equation*}
-u^2 + x^2 + y^2 + z^2 + w^2 = \alpha^2.
\tag{8.4}
\end{equation*}
现在按如下方式在双曲面上诱导坐标 $\{t,\chi,\theta,\phi\}$：
\begin{equation*}
\begin{aligned}
u &= \alpha\sinh(t/\alpha)\\
w &= \alpha\cosh(t/\alpha)\cos\chi\\
x &= \alpha\cosh(t/\alpha)\sin\chi\cos\theta\\
y &= \alpha\cosh(t/\alpha)\sin\chi\sin\theta\cos\phi\\
z &= \alpha\cosh(t/\alpha)\sin\chi\sin\theta\sin\phi.
\end{aligned}
\tag{8.5}
\end{equation*}
% ---- 书页 325 / p340 ----（原书 324/325 分页恰落在本式内：(u,w) 在书页 324，(x,y,z) 及其后文字在书页 325；两块已按原书合并为单个五行公式） ----
于是双曲面上的度规为
\begin{equation*}
\mathrm{d}s^2 = -\mathrm{d}t^2 + \alpha^2\cosh^2(t/\alpha)\left[\mathrm{d}\chi^2 + \sin^2\chi\left(\mathrm{d}\theta^2 + \sin^2\theta\,\mathrm{d}\phi^2\right)\right].
\tag{8.6}
\end{equation*}
我们认出圆括号中的表达式是二维球面上的度规 $\mathrm{d}\Omega_2^2$，而方括号中的表达式是三维球面上的度规 $\mathrm{d}\Omega_3^2$。这样，de Sitter 空间描述的是一个先是收缩、在 $t=0$ 时达到最小尺寸、然后再重新膨胀的空间三维球面。当然，这一特殊描述承袭自某个特定的坐标系；我们将看到，同样有效的其他描述也是存在的。

这些坐标覆盖整个流形。一般你可以这样检验这一点：比如追踪测地线在坐标系边缘附近的行为；如果坐标不完备，测地线看起来就会在有限的仿射参量处终止。因此 de Sitter 空间的拓扑是 $\mathbf{R}\times S^3$。这使得导出共形图变得非常简单，因为构造共形图的关键一步，就是把度规写成它与 Einstein 静态宇宙（一个拓扑为 $\mathbf{R}\times S^3$、描述一个随时间推移保持半径不变的空间三维球面的时空）共形相关联的形式。考虑从 $t$ 到 $t'$ 的坐标变换
\begin{equation*}
\cosh(t/\alpha) = \frac{1}{\cos(t')}.
\tag{8.7}
\end{equation*}
度规 (8.6) 于是变为
\begin{equation*}
\mathrm{d}s^2 = \frac{\alpha^2}{\cos^2(t')}\,\mathrm{d}\bar{s}^2,
\tag{8.8}
\end{equation*}
其中 $\mathrm{d}\bar{s}^2$ 表示 Einstein 静态宇宙上的度规，
\begin{equation*}
\mathrm{d}\bar{s}^2 = -(\mathrm{d}t')^2 + \mathrm{d}\chi^2 + \sin^2\chi\,\mathrm{d}\Omega_2^2.
\tag{8.9}
\end{equation*}
新时间坐标的范围是
\begin{equation*}
-\pi/2 < t' < \pi/2.
\tag{8.10}
\end{equation*}
de Sitter 空间的共形图就是与 de Sitter 共形相关的那一片 Einstein 静态宇宙区域的表示。它看起来像一个正方形，如图 8.1 所示。常数 $t'$ 的类空切片代表一个三维球面；左右两条边缘处的虚线是这个球面的北极和南极。对角线代表类光射线；在过去无限远释放的一个光子，恰好会在未来无限远时到达球面上正对面的那一点（对跖点）。

%FIG{8.1}: {de Sitter 时空的共形图。类空切片是三维球面，因此图上的点代表二维球面，只有左右两条边缘上的点除外——它们代表单个点。}

% ---- 书页 326 / p341 ----
请记住，时空之所以向过去和未来“终结”，仅仅是共形变换的魔力所致；真实的 de Sitter 空间向未来和过去无限延伸。还要注意，两个点可以拥有完全断开的未来（或过去）光锥；这反映了这样一个事实：球状空间截面膨胀得如此之快，以至于来自一点的光永远无法与来自另一点的光相接触。

一个类似的双曲面构造揭示了最大对称而 $\kappa<0$ 的时空，即所谓的 \textbf{反 de Sitter 空间}（anti-de Sitter space）。从一个虚构的五维平直流形出发，其度规为 $\mathrm{d}s_5^2 = -\mathrm{d}u^2 - \mathrm{d}v^2 + \mathrm{d}x^2 + \mathrm{d}y^2 + \mathrm{d}z^2$，并嵌入一个由下式给出的双曲面
\begin{equation*}
-u^2 - v^2 + x^2 + y^2 + z^2 = -\alpha^2.
\tag{8.11}
\end{equation*}
注意这里所有的负号。然后我们可以在双曲面上诱导坐标 $\{t',\rho,\theta,\phi\}$：
\begin{equation*}
\begin{aligned}
u &= \alpha\sin(t')\cosh(\rho)\\
v &= \alpha\cos(t')\cosh(\rho)\\
x &= \alpha\sinh(\rho)\cos\theta\\
y &= \alpha\sinh(\rho)\sin\theta\cos\phi\\
z &= \alpha\sinh(\rho)\sin\theta\sin\phi,
\end{aligned}
\tag{8.12}
\end{equation*}
由此得到这个双曲面上如下形式的度规：
\begin{equation*}
\mathrm{d}s^2 = \alpha^2\left(-\cosh^2(\rho)\,\mathrm{d}t'^2 + \mathrm{d}\rho^2 + \sinh^2(\rho)\,\mathrm{d}\Omega_2^2\right).
\tag{8.13}
\end{equation*}
这些坐标有一个奇怪的特点，即 $t'$ 是周期的。由 (8.12)，$t'$ 与 $t'+2\pi$ 代表双曲面上的同一个位置。由于 $\partial_{t'}$ 处处类时，当 $t'$ 增加时保持 $\{\rho,\theta,\phi\}$ 不变的曲线将是一条闭合类时曲线。然而，这并不是时空的内禀性质，只是我们从特定嵌入导出度规的方式所带来的假象。我们完全可以转而考虑这个流形的“覆盖空间”（covering space），即度规由 (8.13) 给出、但允许 $t'$ 从 $-\infty$ 变到 $\infty$ 的时空。这个空间中没有闭合类时曲线，我们就把这一点取作反 de Sitter 空间的定义。

为了导出共形图，进行一个与 de Sitter 情形类似的坐标变换，不过这次作用在径向坐标上：
\begin{equation*}
\cosh(\rho) = \frac{1}{\cos\chi},
\tag{8.14}
\end{equation*}
于是
\begin{equation*}
\mathrm{d}s^2 = \frac{\alpha^2}{\cos^2\chi}\,\mathrm{d}\bar{s}^2,
\tag{8.15}
\end{equation*}

% ---- 书页 327 / p342 ----
%FIG{8.2}: {反 de Sitter 时空的共形图。类空切片具有 $\mathbf{R}^3$ 的拓扑，我们用极坐标把它表示出来，因此图上的点代表二维球面，只有左边的点除外——它们代表空间原点处的单个点。无限远是右侧的一个类时曲面。}

其中 $\mathrm{d}\bar{s}^2$ 表示 Einstein 静态宇宙的度规 (8.9)。与 de Sitter 不同，径向坐标现在出现在共形因子之中。此外，对反 de Sitter 而言，$t'$ 坐标从负无穷一直到正无穷，而径向坐标的范围是
\begin{equation*}
0 \leq \chi < \frac{\pi}{2}.
\tag{8.16}
\end{equation*}
这样，反 de Sitter 空间与 Einstein 静态宇宙的一半共形相关联。其共形图示于图 8.2，图中画出了几条通过点 $t'=0$、$\chi=0$ 的有代表性的类时和类空测地线。由于 $\chi$ 只到 $\pi/2$ 而不是一直到 $\pi$，这个时空的类空切片具有 $S^3$ 中一个半球面内部的拓扑；也就是说，它在拓扑上是 $\mathbf{R}^3$（因而整个时空具有拓扑 $\mathbf{R}^4$）。注意，我们用极坐标画出这张图：左边的一点代表空间原点处的一个点，而右边的一点代表空间无限远处的一个二维球面。另一种流行的表示方法是画出时空的横截面，让空间原点位于正中间，而右边和左边合起来构成空间无限远。

反 de Sitter 的一个有趣特点是，无限远以一个类时超曲面的形式出现，由 $\chi=\pi/2$ 定义。由于无限远是类时的，这个空间不是全局双曲的，我们就没有一个以类空切片上给定的信息来表述的适定初值问题，因为信息总是可以“从无穷远流入”。另一个有趣的特点是，指数映射并不映满整个时空；如图中所画的那样，从指定点出发的测地线并不覆盖整个流形。指向未来的类时测地线，如图所示，起初可以从 $t'=0$、$\chi=0$ 沿径向向外运动，但最终会重新聚焦到点 $t'=\pi$、$\chi=0$，然后再次沿径向向外运动。

% ---- 书页 328 / p343 ----
顺便说一句，我们忍不住要指出：无限远的类时性质促成了弦论的一个引人注目的特征——“AdS/CFT 对应”（AdS/CFT correspondence）。这里 AdS 当然就是我们一直在讨论的反 de Sitter 空间，而 CFT 指的是定义在边界上的共形不变的场论［对 $n$ 维 AdS 来说，这个边界本身就是一个 $(n-1)$ 维时空］。AdS/CFT 对应表明，在某个极限下，AdS 背景上的量子引力（或其超对称版本）与定义在边界上的共形不变的非引力场论之间存在等价关系。关于非引力的量子场论我们知道得很多，而对量子引力却知之甚少，因此这一对应（如果它是正确的——这看起来很有可能，但尚未被证明）将揭示量子引力中可能发生之事的大量信息。\footnote{全面的综述文章见 O.~Aharony, S.S.~Gubser, J.M.~Maldacena, H.~Ooguri, and Y.~Oz, \emph{Phys.\ Rept.} \textbf{323}, 183 (2000), \texttt{http://arxiv.org/hep-th/9905111}。}

于是我们有了三个最大对称的时空：Minkowski（$\kappa=0$）、de Sitter（$\kappa>0$）和反 de Sitter（$\kappa<0$）。它们中有哪一个是真实世界的有用模型吗？进而言之，它们是 Einstein 方程的解吗？先对 (8.1) 给出的 Riemann 张量取迹，并限定到四维：
\begin{equation*}
R_{\mu\nu} = 3\kappa g_{\mu\nu},\qquad R = 12\kappa.
\tag{8.17}
\end{equation*}
可见在最大对称空间中，Ricci 张量与度规成正比。具有这种性质的时空有时被称为 Einstein 空间（Einstein space）；而 Einstein 静态宇宙并\emph{不}是 Einstein 空间的一个例子，这一点有时会令人混淆。更糟的是，我们后面还会遇到 Einstein–de Sitter 宇宙学，它既与 Einstein 空间无关，也与 Einstein 静态宇宙或 de Sitter 空间无关。Einstein 张量为
\begin{equation*}
G_{\mu\nu} = R_{\mu\nu} - \frac{1}{2}Rg_{\mu\nu} = -3\kappa g_{\mu\nu}.
\tag{8.18}
\end{equation*}
因此，Einstein 方程 $G_{\mu\nu}=8\pi G T_{\mu\nu}$ 蕴含（在最大对称时空中如此，并非普遍情形）能量-动量张量与度规成正比：
\begin{equation*}
T_{\mu\nu} = -\frac{3\kappa}{8\pi G}\,g_{\mu\nu}.
\tag{8.19}
\end{equation*}
这样的能量-动量张量对应于真空能或宇宙学常数，如第 4 章所讨论的。能量密度和压强为
\begin{equation*}
\rho = -p = \frac{3\kappa}{8\pi G}.
\tag{8.20}
\end{equation*}
若 $\rho$ 为正，我们得到 de Sitter 解；若 $\rho$ 为负，我们得到反 de Sitter。但在我们的宇宙中，除了可能有真空能之外，还有普通的物质和辐射。我们的最大对称时空并不相容于
% ---- 书页 329 / p344 ----
动力学上可观的物质和/或辐射。此外，由于我们观测到宇宙中的可见物质正在彼此分开（宇宙正在膨胀，下面会讨论），物质的密度在过去更高；所以即使物质对总能量的贡献在今天可以忽略，它在更早的宇宙中也曾经是可观的。因此，最大对称时空并不是真实世界的合理模型。不过，它们确实代表了在没有任何普通物质或引力辐射时 Einstein 方程的（局域）唯一解；正是在这个意义上，它们可以被看作广义相对论的基态。

\section{Robertson–Walker 度规}

为了描述真实世界，我们不得不放弃“完美的”哥白尼原理——它蕴含空间和时间处处都对称——而假定某种更宽容的东西。事实证明，直接假定宇宙在\emph{空间上}均匀且各向同性、但随时间演化，既是直截了当的，又与观测一致。在广义相对论中，这转译为如下陈述：宇宙可以叶状分层（foliated）为一些类空切片，使得每个三维切片都是最大对称的。于是我们把时空取为 $\mathbf{R}\times\Sigma$，其中 $\mathbf{R}$ 代表时间方向，而 $\Sigma$ 是一个最大对称的三维流形。时空度规因此取如下形式
\begin{equation*}
\mathrm{d}s^2 = -\mathrm{d}t^2 + R^2(t)\,\mathrm{d}\sigma^2,
\tag{8.21}
\end{equation*}
其中 $t$ 是类时坐标，$R(t)$ 是一个称为\textbf{尺度因子}（scale factor）的函数，而 $\mathrm{d}\sigma^2$ 是 $\Sigma$ 上的度规，它可以表示为
\begin{equation*}
\mathrm{d}\sigma^2 = \gamma_{ij}(u)\,\mathrm{d}u^i\,\mathrm{d}u^j,
\tag{8.22}
\end{equation*}
其中 $(u^1,u^2,u^3)$ 是 $\Sigma$ 上的坐标，$\gamma_{ij}$ 是一个最大对称的三维度规。尺度因子告诉我们类空切片 $\Sigma$ 在时刻 $t$ 有多大。（不要把它与标量曲率混淆。）这里使用的坐标，其度规没有交叉项 $\mathrm{d}t\,\mathrm{d}u^i$，且 $\mathrm{d}t^2$ 的系数与 $u^i$ 无关，称为\textbf{共动坐标}（comoving coordinates），它是附录 D 讨论的高斯法坐标的一个特例。保持在常数 $u^i$ 处的观者也称为“共动的”。只有共动观者才会认为宇宙看起来是各向同性的；事实上在地球上我们并不完全共动，结果便是我们在宇宙微波背景中看到偶极各向异性，它来源于常规的多普勒效应。

因此我们感兴趣的是最大对称的欧几里得三维度规 $\gamma_{ij}$。我们知道，最大对称的度规满足
\begin{equation*}
{}^{(3)}R_{ijkl} = k\left(\gamma_{ik}\gamma_{jl} - \gamma_{il}\gamma_{jk}\right),
\tag{8.23}
\end{equation*}
其中为了今后的方便我们引入了
\begin{equation*}
k = {}^{(3)}R/6,
\tag{8.24}
\end{equation*}

% ---- 书页 330 / p345 ----
并且在 Riemann 张量上放一个上标 $(3)$，提醒我们它与三维度规 $\gamma_{ij}$ 相联系，而不是与整个时空的度规相联系。Ricci 张量便是
\begin{equation*}
{}^{(3)}R_{jl} = 2k\gamma_{jl}.
\tag{8.25}
\end{equation*}
如果空间要是最大对称的，那它肯定是球对称的。从对 Schwarzschild 解的探索中，我们已经对球对称空间有所了解；度规可以写成
\begin{equation*}
\mathrm{d}\sigma^2 = \gamma_{ij}\,\mathrm{d}u^i\,\mathrm{d}u^j = e^{2\beta(\bar{r})}\,\mathrm{d}\bar{r}^2 + \bar{r}^2\,\mathrm{d}\Omega^2,
\tag{8.26}
\end{equation*}
其中 $\bar{r}$ 是径向坐标，二维球面上的度规一如既往是 $\mathrm{d}\Omega^2 = \mathrm{d}\theta^2 + \sin^2\theta\,\mathrm{d}\phi^2$。这种度规的 Ricci 张量分量可以从 (5.14)——静态球对称时空的 Ricci 张量——通过令 $\alpha=0$ 和 $r=\bar{r}$ 得到，即
\begin{equation*}
\begin{aligned}
{}^{(3)}R_{11} &= \frac{2}{\bar{r}}\,\partial_1\beta\\
{}^{(3)}R_{22} &= e^{-2\beta}\left(\bar{r}\,\partial_1\beta - 1\right) + 1\\
{}^{(3)}R_{33} &= \left[e^{-2\beta}\left(\bar{r}\,\partial_1\beta - 1\right) + 1\right]\sin^2\theta.
\end{aligned}
\tag{8.27}
\end{equation*}
我们利用 (8.25) 令它们与度规成正比，就可以解出 $\beta(\bar{r})$：
\begin{equation*}
\beta = -\frac{1}{2}\ln\left(1-k\bar{r}^2\right),
\tag{8.28}
\end{equation*}
由此得到三维曲面 $\Sigma$ 上的度规
\begin{equation*}
\mathrm{d}\sigma^2 = \frac{\mathrm{d}\bar{r}^2}{1-k\bar{r}^2} + \bar{r}^2\,\mathrm{d}\Omega^2.
\tag{8.29}
\end{equation*}
从 (8.24) 可以注意到，$k$ 的值决定了空间曲面的曲率，从而决定了空间曲面的大小。通常把它归一化，使得
\begin{equation*}
k \in \{+1, 0, -1\},
\tag{8.30}
\end{equation*}
而把流形的物理尺寸吸收进尺度因子 $R(t)$。

$k=-1$ 的情形对应于 $\Sigma$ 上曲率为负的常数，有时称为\textbf{开放}（open）；$k=0$ 的情形对应于 $\Sigma$ 上没有曲率，称为\textbf{平直}（flat）；$k=+1$ 的情形对应于 $\Sigma$ 上曲率为正，有时称为\textbf{闭合}（closed）。引入一个由下式定义的新径向坐标 $\chi$，得到度规的另一种形式，可以使这些情形的物理解释更加清楚：
\begin{equation*}
\mathrm{d}\chi = \frac{\mathrm{d}\bar{r}}{\sqrt{1-k\bar{r}^2}}.
\tag{8.31}
\end{equation*}

% ---- 书页 331 / p346 ----
这可以积分得到
\begin{equation*}
\bar{r} = S_k(\chi),
\tag{8.32}
\end{equation*}
其中
\begin{equation*}
S_k(\chi) \equiv
\begin{cases}
\sin(\chi), & k = +1\\
\chi, & k = 0\\
\sinh(\chi), & k = -1,
\end{cases}
\tag{8.33}
\end{equation*}
于是
\begin{equation*}
\mathrm{d}\sigma^2 = \mathrm{d}\chi^2 + {S_k}^2(\chi)\,\mathrm{d}\Omega^2.
\tag{8.34}
\end{equation*}
对平直情形 $k=0$，$\Sigma$ 上的度规变为
\begin{equation*}
\begin{aligned}
\mathrm{d}\sigma^2 &= \mathrm{d}\chi^2 + \chi^2\,\mathrm{d}\Omega^2\\
&= \mathrm{d}x^2 + \mathrm{d}y^2 + \mathrm{d}z^2,
\end{aligned}
\tag{8.35}
\end{equation*}
这就是平直的欧几里得空间。整体上，它可以描述 $\mathbf{R}^3$ 或者更复杂的流形，比如三维环面 $S^1\times S^1\times S^1$。对闭合情形 $k=+1$，我们有
\begin{equation*}
\mathrm{d}\sigma^2 = \mathrm{d}\chi^2 + \sin^2\chi\,\mathrm{d}\Omega^2,
\tag{8.36}
\end{equation*}
这是三维球面的度规。在这种情形下，唯一可能的整体结构是完全的三维球面（不可定向流形 $\mathbf{RP}^3$ 除外，它由认同 $S^3$ 上的对跖点得到）。最后，对开放的 $k=-1$ 情形，我们得到
\begin{equation*}
\mathrm{d}\sigma^2 = \mathrm{d}\chi^2 + \sinh^2\chi\,\mathrm{d}\Omega^2.
\tag{8.37}
\end{equation*}
这是常负曲率三维空间的度规，是 3.9 节讨论的那个双曲面的推广。整体上，这样的空间可以无限延伸（这正是“开放”一词的由来），但它也可以描述一个非单连通的紧致空间（所以“开放”实在算不上最准确的描述）。

时空的度规描述的正是这些最大对称超曲面之一在尺寸上的演化，它可以写成
\begin{equation*}
\mathrm{d}s^2 = -\mathrm{d}t^2 + R^2(t)\left[\frac{\mathrm{d}\bar{r}^2}{1-k\bar{r}^2} + \bar{r}^2\,\mathrm{d}\Omega^2\right].
\tag{8.38}
\end{equation*}
这就是 \textbf{Robertson–Walker（RW）度规}。我们还没有用到 Einstein 方程；它将决定尺度因子 $R(t)$ 的行为。注意，代换

% ---- 书页 332 / p347 ----
\begin{equation*}
\begin{aligned}
R &\to \lambda^{-1}R\\
\bar{r} &\to \lambda\bar{r}\\
k &\to \lambda^{-2}k
\end{aligned}
\tag{8.39}
\end{equation*}
使 (8.38) 保持不变。因此我们可以选取方便的归一化。在曲率 $k$ 归一化为 $\{+1,0,-1\}$ 的那些变量中，尺度因子具有距离的量纲，而径向坐标 $\bar{r}$（或 $\chi$）实际上是无量纲的；这是最流行的选择。我们偏要反其道而行之，改用无量纲的尺度因子
\begin{equation*}
a(t) = \frac{R(t)}{R_0},
\tag{8.40}
\end{equation*}
一个具有距离量纲的坐标
\begin{equation*}
r = R_0\bar{r},
\tag{8.41}
\end{equation*}
以及一个具有 (长度)$^{-2}$ 量纲的曲率参数
\begin{equation*}
\kappa = \frac{k}{R_0^2}.
\tag{8.42}
\end{equation*}
注意，$\kappa$ 可以取任意值，而不仅仅是 $\{+1,0,-1\}$。在这些变量中，Robertson–Walker 度规为
\begin{equation*}
\boxed{\;\mathrm{d}s^2 = -\mathrm{d}t^2 + a^2(t)\left[\frac{\mathrm{d}r^2}{1-\kappa r^2} + r^2\,\mathrm{d}\Omega^2\right].\;}
\tag{8.43}
\end{equation*}
要转换到更常见的记号，只需代入关系式 (8.40)、(8.41) 和 (8.42)。

有了度规在手，我们就可以着手计算联络系数和曲率张量。令 $\dot{a}\equiv \mathrm{d}a/\mathrm{d}t$，Christoffel 符号为
\begin{equation*}
\begin{array}{ll}
\displaystyle \Gamma^0_{11} = \frac{a\dot{a}}{1-\kappa r^2} & \displaystyle \Gamma^1_{11} = \frac{\kappa r}{1-\kappa r^2}\\[10pt]
\displaystyle \Gamma^0_{22} = a\dot{a}r^2 & \displaystyle \Gamma^0_{33} = a\dot{a}r^2\sin^2\theta\\[6pt]
\displaystyle \Gamma^1_{01} = \Gamma^2_{02} & \displaystyle \Gamma^3_{03} = \frac{\dot{a}}{a}\\[6pt]
\displaystyle \Gamma^1_{22} = -r\left(1-\kappa r^2\right) & \displaystyle \Gamma^1_{33} = -r\left(1-\kappa r^2\right)\sin^2\theta\\[6pt]
\displaystyle \Gamma^2_{12} = \Gamma^3_{13} = \frac{1}{r} & \\[6pt]
\displaystyle \Gamma^2_{33} = -\sin\theta\cos\theta & \displaystyle \Gamma^3_{23} = \cot\theta,
\end{array}
\tag{8.44}
\end{equation*}

% ---- 书页 333 / p348 ----
凡未明确写出的，或者为零，或者由对称性与这些相联系。Ricci 张量的非零分量为
\begin{equation*}
\begin{aligned}
R_{00} &= -3\frac{\ddot{a}}{a}\\[4pt]
R_{11} &= \frac{a\ddot{a} + 2\dot{a}^2 + 2\kappa}{1-\kappa r^2}\\[4pt]
R_{22} &= r^2\left(a\ddot{a} + 2\dot{a}^2 + 2\kappa\right)\\
R_{33} &= r^2\left(a\ddot{a} + 2\dot{a}^2 + 2\kappa\right)\sin^2\theta,
\end{aligned}
\tag{8.45}
\end{equation*}
而 Ricci 标量则为
\begin{equation*}
R = 6\left[\frac{\ddot{a}}{a} + \left(\frac{\dot{a}}{a}\right)^2 + \frac{\kappa}{a^2}\right].
\tag{8.46}
\end{equation*}

\section{弗里德曼方程}

RW 度规对尺度因子 $a(t)$ 的任何行为都有定义；我们的下一步将是把它代入 Einstein 方程，导出把尺度因子与宇宙的能量-动量联系起来的弗里德曼方程（组）（Friedmann equation(s)）。我们将选择用理想流体来模型化物质和能量。显然，如果某种流体在某一个参考系中是各向同性的，而它导致的度规在某一个参考系中是各向同性的，那么这两个参考系将会重合；也就是说，流体在共动坐标中静止。四速度就是
\begin{equation*}
U^\mu = (1, 0, 0, 0),
\tag{8.47}
\end{equation*}
而能量-动量张量
\begin{equation*}
T_{\mu\nu} = (\rho + p)U_\mu U_\nu + pg_{\mu\nu}
\tag{8.48}
\end{equation*}
变为
\begin{equation*}
T_{\mu\nu} = \left(\begin{array}{cccc}
\rho & 0 & 0 & 0\\
0 & & & \\
0 & & g_{ij}p & \\
0 & & &
\end{array}\right).
\tag{8.49}
\end{equation*}
升起一个指标后，它取如下方便的形式：
\begin{equation*}
T^\mu{}_\nu = \mathrm{diag}(-\rho, p, p, p).
\tag{8.50}
\end{equation*}
注意，迹由下式给出：
\begin{equation*}
T = T^\mu{}_\mu = -\rho + 3p.
\tag{8.51}
\end{equation*}
在代入 Einstein 方程之前，先考察能量守恒方程的零分量是有教益的：

% ---- 书页 334 / p349 ----
\begin{equation*}
\begin{aligned}
0 &= \nabla_\mu T^\mu{}_{0}\\
&= \partial_\mu T^\mu{}_{0} + \Gamma^\mu_{\mu\lambda}T^\lambda{}_{0} - \Gamma^\lambda_{\mu 0}T^\mu{}_{\lambda}\\
&= -\partial_0\rho - 3\frac{\dot{a}}{a}\left(\rho + p\right).
\end{aligned}
\tag{8.52}
\end{equation*}
为了取得进展，我们可以选取一个\textbf{物态方程}（equation of state），即 $\rho$ 与 $p$ 之间的一个关系。与宇宙学相关的理想流体常常服从简单的物态方程
\begin{equation*}
p = w\rho,
\tag{8.53}
\end{equation*}
其中 $w$ 是与时间无关的常数。当然，无论它是否保持为常数，我们都可以随意定义参数 $w=p/\rho$；然而，如果 $w$ 在变化，把 $p=w\rho$ 称为“物态方程”就不那么名副其实了。能量守恒方程变为
\begin{equation*}
\boxed{\;\frac{\dot{\rho}}{\rho} = -3(1+w)\frac{\dot{a}}{a}.\;}
\tag{8.54}
\end{equation*}
如果 $w$ 是常数，这可以积分得到
\begin{equation*}
\rho \propto a^{-3(1+w)}.
\tag{8.55}
\end{equation*}
要想了解 $w$ 允许取什么样的值，请参考第 4 章关于能量条件的讨论。类光主能量条件（Null Dominant Energy Condition）允许任何符号的真空能，但除此之外要求物质不能使真空失稳，它蕴含
\begin{equation*}
|w| \leq 1.
\tag{8.56}
\end{equation*}
这个要求绝非不可动摇的金科玉律，但作为研究真实世界中可能发生之事的出发点，它看起来是明智而保守的。

宇宙学流体最流行的两个例子称为\textbf{物质}（matter）和\textbf{辐射}（radiation）。物质是任何一组无碰撞、非相对论性的粒子，其压强基本上为零：
\begin{equation*}
p_{\mathrm{M}} = 0.
\tag{8.57}
\end{equation*}
例子包括普通的恒星和星系，对它们来说，压强与能量密度相比可以忽略。物质也称为尘埃（dust），能量密度主要由物质构成的宇宙称为\textbf{物质主导}的（matter-dominated）。物质中的能量密度按
\begin{equation*}
\rho_{\mathrm{M}} \propto a^{-3}.
\tag{8.58}
\end{equation*}
衰减。这可以直接解释为：随着宇宙膨胀，粒子数密度在减少。对物质而言，能量密度由静能主导，
% ---- 书页 335 / p350 ----
而静能与数密度成正比。辐射既可以描述真实的电磁辐射，也可以描述相对速度足够接近光速、以至于（至少就其物态方程而言）与光子无法区分的有质量粒子。虽然各向同性的相对论性粒子气体是一种理想流体，因而具有由 (8.48) 给出的能量-动量张量，我们也知道，电磁学的 $T_{\mu\nu}$ 可以用场强表达为
\begin{equation*}
T^{\mu\nu} = F^{\mu\lambda}F^{\nu}{}_{\lambda} - \frac{1}{4}g^{\mu\nu}F^{\lambda\sigma}F_{\lambda\sigma}.
\tag{8.59}
\end{equation*}
它的迹由下式给出：
\begin{equation*}
T^\mu{}_\mu = F^{\mu\lambda}F_{\mu\lambda} - \frac{1}{4}(4)F^{\lambda\sigma}F_{\lambda\sigma} = 0.
\tag{8.60}
\end{equation*}
但它也必须等于 (8.51)，所以物态方程为
\begin{equation*}
p_{\mathrm{R}} = \frac{1}{3}\rho_{\mathrm{R}}.
\tag{8.61}
\end{equation*}
能量密度大部分以辐射形式存在的宇宙称为\textbf{辐射主导}的（radiation-dominated）。辐射中的能量密度按
\begin{equation*}
\rho_{\mathrm{R}} \propto a^{-4}.
\tag{8.62}
\end{equation*}
衰减。这样，辐射中的能量密度比物质中的衰减得稍快一些；这是因为光子数密度的减少方式与非相对论性粒子数密度的减少方式相同，但单个光子在红移时还会以 $a^{-1}$ 的方式损失能量，这一点我们稍后会看到。同样，有质量但相对论性的粒子在共动坐标中“慢下来”时也会损失能量。我们相信，今天辐射的能量密度远小于物质的能量密度，$\rho_{\mathrm{M}}/\rho_{\mathrm{R}}\sim 10^{3}$。然而在过去，宇宙要小得多，在极早的时刻，辐射的能量密度本会占据主导。

正如我们讨论过的，真空能也取理想流体的形式，其物态方程为 $p_{\Lambda}=-\rho_{\Lambda}$。其能量密度为常数，
\begin{equation*}
\rho_{\Lambda} \propto a^{0}.
\tag{8.63}
\end{equation*}
由于物质和辐射的能量密度随着宇宙膨胀而减少，如果存在非零的真空能，那么只要宇宙不开始收缩，从长远来看它往往会占上风。如果发生这种情况，我们就说宇宙变得\textbf{真空主导}（vacuum-dominated）。de Sitter 空间和反 de Sitter 空间都是真空主导的解。

现在我们转向 Einstein 方程。回想它可以写成 (4.45) 的形式：
\begin{equation*}
R_{\mu\nu} = 8\pi G\left(T_{\mu\nu} - \frac{1}{2}g_{\mu\nu}T\right).
\tag{8.64}
\end{equation*}
